\chapter{Technologie wykorzystane w projekcie}
\label{cha:technologie}

\section{Technologie interfejsu użytkownika}
\label{sec:techUI}

Warstwę prezentacji zrealizowano w technologii \emph{Windows Presentation Foundation}
(WPF)~\cite{msdocsWPF} na platformie .NET, w języku C\#~\cite{msdocsCSharp}. WPF posłużył do budowy
okien aplikacji, paska bocznego nawigacji oraz ekranów poszczególnych modułów. Wygląd aplikacji
(kolory, style kontrolek) zdefiniowano w słowniku zasobów \texttt{Theme.xaml}, co oddziela warstwę
wizualną od logiki. Strona logowania zbudowana jest deklaratywnie w pliku XAML, natomiast pozostałe
ekrany tworzone są dynamicznie w kodzie, z wykorzystaniem wspólnych metod pomocniczych budujących
powtarzalne elementy interfejsu (karty, znaczniki statusu, wykresy).

\section{Technologie dostępu do danych}
\label{sec:techDane}

Warstwę dostępu do danych oparto na bibliotece \emph{Entity Framework Core}~\cite{msdocsEFCore} -
mechanizmie mapowania obiektowo-relacyjnego (ORM). Pozwala on operować na bazie za pomocą obiektów
języka C\# i zapytań LINQ, bez ręcznego pisania większości zapytań SQL. Systemem zarządzania bazą danych jest
\emph{PostgreSQL}~\cite{postgresql}, łączony z aplikacją za pośrednictwem dostawcy
\emph{Npgsql}~\cite{npgsql}. Schemat bazy jest wersjonowany i tworzony automatycznie poprzez migracje
Entity Framework Core.

\section{Dodatkowe biblioteki i narzędzia}
\label{sec:techDodatkowe}

W projekcie wykorzystano ponadto:

\begin{itemize}
  \item \texttt{System.Security.Cryptography} - moduł platformy .NET użyty do wyznaczania skrótu
        SHA-256 haseł użytkowników.
  \item \texttt{System.IO} (\texttt{StreamWriter}) - strumieniowy zapis pliku CSV z wynikami testu,
        z kodowaniem UTF-8 obsługującym polskie znaki.
  \item \emph{Visual Studio} - środowisko programistyczne wykorzystane do budowy i uruchamiania
        aplikacji.
  \item \emph{pgAdmin} - narzędzie do zarządzania bazą danych PostgreSQL.
  \item \emph{Git} - system kontroli wersji projektu.
\end{itemize}

\section{Uzasadnienie wyboru technologii}
\label{sec:techUzasadnienie}

WPF, Entity Framework Core i PostgreSQL dobrano pod charakter aplikacji:

\begin{itemize}
  \item \textbf{WPF} wybrano, bo aplikacja jest programem desktopowym z rozbudowanym interfejsem
        (karty, wykresy, stylizacja), a nie potrzebuje serwera ani przeglądarki.
  \item \textbf{Entity Framework Core} wybrano zamiast bezpośredniego ADO.NET, ponieważ model danych
        opiera się na dziedziczeniu (trzy hierarchie klas) i wielu relacjach. ORM odwzorowuje je
        automatycznie (strategią Table-per-Hierarchy), a migracje pozwalają wersjonować schemat bez
        ręcznego pisania skryptów SQL, co skraca kod warstwy danych.
  \item \textbf{PostgreSQL} to darmowy, otwarty i niezawodny system bazodanowy, dobrze wspierany
        przez Entity Framework Core poprzez dostawcę Npgsql.
\end{itemize}
